\documentclass[a4paper,10pt]{article}

\usepackage[utf8]{inputenc}
\usepackage[T1]{fontenc}
\usepackage[brazil]{babel}
\usepackage{enumitem}
\usepackage{hyperref}
\usepackage{geometry}
\geometry{margin=1in}

\hypersetup{
    pdftitle={Currículo - Gian Pedro Rodrigues},
    pdfauthor={Gian Pedro Rodrigues},
    pdfsubject={Currículo Profissional},
    colorlinks=true,
    urlcolor=blue
}

\setlength{\parindent}{0pt}
\setlength{\parskip}{6pt}

\begin{document}

\vspace*{-1cm}
\begin{center}
    {\LARGE \textbf{Gian Pedro Rodrigues}}\\
    \vspace{0.2cm}
    \href{mailto:gianpedrodev@gmail.com}{gianpedrodev@gmail.com} \quad | \quad (12) 99618-3274\\
    \href{https://github.com/GianSE}{github.com/GianSE} \quad | \quad
    \href{https://www.linkedin.com/in/gian-pedro-rodrigues-3b6a44259/}{linkedin.com/in/gian-pedro-rodrigues}
\end{center}

\vspace{0.4cm}

\section*{Resumo Profissional}
Estudante de Engenharia de Computação na UTFPR, com experiência prática em \textbf{Análise de Dados}, atuando na Drogamais com \textbf{Python, SQL e Automação} para otimização de processos e análise de informações. Possui também base em desenvolvimento Full Stack (Node.js, React, Next.js) e projetos em Machine Learning. Atua em projetos reais através de empresa júnior, iniciativas acadêmicas e extracurriculares.

\section*{Objetivos}
Desenvolver minha carreira profissional atuando na área de \textbf{Dados}, aplicando conhecimentos em Python, SQL e Automação. E também desenvolver minhas habilidades em desenvolvimento de \textbf{software/web} com ReactJS, NodeJS e Python.

\section*{Formação Acadêmica}
\begin{itemize}[leftmargin=*]
    \item \textbf{UTFPR – Universidade Tecnológica Federal do Paraná} \hfill \textit{Jul 2022 – Jun 2027 (em andamento)}\\
    Bacharelado em Engenharia de Computação

    \item \textbf{ETEC – Escola Técnica Estadual} \hfill \textit{Jan 2015 – Dez 2017}\\
    Técnico em Contabilidade
\end{itemize}

\section*{Experiências}
\begin{itemize}[leftmargin=*]
    \item \textbf{Analista de Dados – Drogamais} \hfill \textit{2025 – atual}\\
    Responsável pela criação de scripts para a extração de dados por API ou por meio de automação como Power Automate ou Selenium, processamento e tranformação de dados através das linguagens \textbf{Python}(pandas) e \textbf{SQL}. E criação de data view para diretoria e lojistas, Power BI, para contribuir para decisões estratégicas para o varejo farmacêutico e eficiência operacional.

    \item \textbf{Desenvolvedor Back-End – Empresa Júnior Unect} \hfill \textit{2024 – 2025}\\
    Atuação com Node.js, Express.js, JWT, React.js, Next.js e bancos de dados (MongoDB) no desenvolvimento de soluções para aplicações Web, Desktop e Mobile.

    \item \textbf{Desenvolvedor – Projeto LtSat UTFPR} \hfill \textit{2024 – 2025}\\
    Desenvolveu sistema de recepção e visualização de dados de um satélite (STM32), com back-end em Node.js e exibição interativa em D3.js no front-end. Participou da CanSat Competition nos EUA, representando o Brasil como a única equipe brasileira.
\end{itemize}

\section*{Habilidades}
\begin{itemize}[leftmargin=*]
    \item \textbf{Dados e Automação:} Python, SQL (PostgreSQL, MariaDB, MongoDB), Pandas, NumPy, Automação, TensorFlow, Keras, Scikit-learn
    \item \textbf{Back-End:} Node.js, Express.js, JWT, Java Spring Boot
    \item \textbf{Front-End:} React.js, Next.js, HTML, CSS, Tailwind, Bootstrap
    \item \textbf{Ferramentas:} Git, GitHub, Linux, WSL, Postman, QtCreator
    \item \textbf{Outras Linguagens:} C, C++
\end{itemize}

\section*{Idiomas}
\begin{itemize}[leftmargin=*]
    \item \textbf{Inglês:} Intermediário (B1) – leitura, escrita, escuta e conversação. Formação pelo CALEM (Centro Acadêmico de Línguas Estrangeiras Modernas) na UTFPR.
\end{itemize}

\section*{Projetos Educacionais de Software}
\begin{itemize}[leftmargin=*]
    \item \textbf{API para Clínica Veterinária}\\
    Desenvolveu uma API RESTful completa com Node.js e Express, incluindo autenticação via JWT e integração com MongoDB Atlas. Utilização do Postman para testes automatizados e organização de fluxo de requisições com scripts.\\
    \href{https://github.com/GianSE/capacitacao-parte2}{github.com/GianSE/capacitacao-parte2}

    \item \textbf{UTF em Foco}\\
    Desenvolveu um site informativo responsivo como projeto acadêmico, utilizando HTML, CSS (Bootstrap) e JavaScript. Estrutura modular e layout adaptável para diferentes dispositivos.\\
    \href{https://gianse.github.io/UTF_em_Foco/}{gianse.github.io/UTF_em_Foco/}

    \item \textbf{ClimaTech (Desktop)}\\
    Desenvolveu aplicação desktop em C++ (Qt Creator) com integração à API de clima (OpenWeather) e banco de dados PostgreSQL para persistência dos dados consultados.\\
    \href{https://github.com/GianSE/poo2_climatech}{github.com/GianSE/poo2_climatech}

    \item \textbf{MobileNet – Cifar10}\\
    Desenvolveu um modelo de rede neural convolucional usando MobileNet para classificação de imagens no conjunto de dados CIFAR-10. O projeto utilizou TensorFlow, Keras e técnicas de pré-processamento e avaliação de desempenho em visão computacional.\\
    \href{https://colab.research.google.com/drive/1AGP831WCSEpkeN7zi-tSyno52o2E8LfE}{Notebook no Google Colab do treinamento}

    \item \textbf{Simulação de Tunelamento Quântico (Iniciação Científica)}\\
    Desenvolveu uma simulação do fenômeno de tunelamento quântico utilizando Python. Modelagem matemática da equação de Schrödinger em potenciais unidimensionais, com visualizações gráficas usando Matplotlib e geração de animações com FFmpeg. Projeto voltado à compreensão visual de fenômenos da mecânica quântica.\\
    \href{https://colab.research.google.com/drive/122zQ2S9tb0PJ_tqczK8SwqvF77C_w4Md?usp=sharing}{Notebook no Google Colab da simulação}
\end{itemize}

\end{document}
